\documentclass[11pt]{article}
\usepackage[T1]{fontenc}
\usepackage{lmodern}
\usepackage{amsmath,amssymb}
\usepackage{booktabs}
\usepackage{array}
\usepackage[a4paper,margin=2.5cm]{geometry}
\usepackage[hidelinks]{hyperref}
\newcommand{\doi}[1]{\href{https://doi.org/#1}{\texttt{#1}}}
\setlength{\parskip}{4pt}
\title{\vspace{-1.6cm}\Large Piezoelectric Semiconductor Plate: Reduced Field Equations, Solution Form\\
and Characteristic Matrix\\[8pt]\large Companion to the research plan, third file}
\author{Dr.\ Vipin \\ Gurugram University}
\date{30 October 2026}
\begin{document}
\maketitle
\vspace{-6mm}

\section{What this file is}

The research plan stops at the field equations. This file goes one step further and writes down the
plane-wave reduction of those equations, the form of the solution through the thickness, and the matrix whose
determinant gives the dispersion relation. Only results are written. None of the intermediate algebra is
shown. You are expected to redo it on paper, term by term, and to find the places where this file or the
sources are wrong. Every formula carries a source or says that it is ours. A formula marked ``ours'' has no
paper behind it, only a law you can look up, and it needs the most careful look.

Equation numbers refer to this file. Equations (1) to (7) of the plan are called Plan (1) to Plan (7).
The reference list is the one in the research plan, plus the entries listed at the end of this file.

\section{Setting}

The plate occupies $-h\le z\le h$ and extends without limit in $x$ and $y$. The crystal is hexagonal
(class 6mm) with its $c$-axis along $z$. We take plane strain, $u_y=0$, $\partial_y=0$, and one plane wave in
$x$,
\begin{equation}
\bigl(u_x,u_z,\varphi,\Delta n,\theta\bigr)
=\bigl(\hat u_x,\hat u_z,\hat\varphi,\hat n,\hat\theta\bigr)(z)\;e^{\,i(kx-\omega t)} ,
\label{eq:ansatz}
\end{equation}
with $k=k_r+ik_i$ complex and $\omega$ real. Through the thickness every field is a sum of exponentials
$e^{\eta z}$. The constants that survive in the 6mm class with this geometry are $c_{11},c_{13},c_{33},c_{44}$,
$e_{31},e_{33},e_{15}$, $\epsilon_{11},\epsilon_{33}$, the thermal stress coefficients $\beta_1,\beta_3$
(the two independent components of $\beta_{ij}=c_{ijkl}\alpha_{kl}$) and the pyroelectric constant $p_3$.
The tensor-to-matrix rule is the usual one: $xx\to1$, $zz\to3$, $xz\to5$. Please write out which elements
of $c_{ijkl}$, $e_{kij}$ and $\epsilon_{ik}$ you use and why the others vanish; this is part of your work.

\section{The stresses and the displacement in component form}

Put $\nabla^2\to\eta^2-k^2$ and define two operator factors,
\begin{equation}
\mathcal N(\eta)=1+l_{\rm nl}^2\,(k^2-\eta^2),\qquad
\mathcal G(\eta)=1+l_g^2\,(k^2-\eta^2).
\label{eq:NG}
\end{equation}
They are Plan (1) with the Laplacian acting on $e^{\eta z}e^{ikx}$. With $\partial_x\to ik$, $\partial_z\to\eta$,
the strains and the field are $\varepsilon_{xx}=ik\hat u_x$, $\varepsilon_{zz}=\eta\hat u_z$,
$\varepsilon_{xz}=\tfrac12(\eta\hat u_x+ik\hat u_z)$, $E_x=-ik\hat\varphi$, $E_z=-\eta\hat\varphi$
\cite[Eq.\ (2)]{Zhu2023}. The three stresses and the two components of the electric displacement of Plan (1)--(2) are then
\begin{align}
\hat\sigma_{xx}&=\mathcal G\,(ik\,c_{11}\hat u_x+\eta\,c_{13}\hat u_z)+\eta\,e_{31}\hat\varphi-\beta_1\hat\theta,\\
\hat\sigma_{zz}&=\mathcal G\,(ik\,c_{13}\hat u_x+\eta\,c_{33}\hat u_z)+\eta\,e_{33}\hat\varphi-\beta_3\hat\theta,\\
\hat\sigma_{xz}&=\mathcal G\,c_{44}\,(\eta\hat u_x+ik\hat u_z)+ik\,e_{15}\hat\varphi,\\
\hat D_x&=e_{15}(\eta\hat u_x+ik\hat u_z)-ik\,\epsilon_{11}\hat\varphi,\\
\hat D_z&=ik\,e_{31}\hat u_x+\eta\,e_{33}\hat u_z-\eta\,\epsilon_{33}\hat\varphi+p_3\hat\theta .
\end{align}
Source: the structure is Plan (1)--(2), which follows \cite[Eq.\ (1)]{Zhu2023} for the piezoelectric part and
\cite[Eq.\ (3.1)]{Quintanilla2019} for the thermal stress; the component form is ours.

\section{The five reduced equations and the matrix $\mathbf M(\eta)$}

The five unknown amplitudes are collected in $\mathbf a=(\hat u_x,\hat u_z,\hat\varphi,\hat n,\hat\theta)^{T}$.
Plan (3) to Plan (6), with the substitutions above, become $\mathbf M(\eta)\,\mathbf a=\mathbf 0$ with
$L=\eta^2-k^2$ and $f=1-i\omega\tau_q$. The sources of the five equations are given in the plan. Row by row,
the entries of $\mathbf M$ are the following. Entries not listed are zero.

\paragraph{Momentum in $x$ (Plan (3), $x$-component).}
\begin{align}
M_{11}&=-\mathcal N\rho\omega^2+\mathcal G\,(c_{11}k^2-c_{44}\eta^2), &
M_{12}&=-ik\eta\,\mathcal G\,(c_{13}+c_{44}),\\
M_{13}&=-ik\eta\,(e_{31}+e_{15}), &
M_{15}&=ik\,\beta_1 .
\end{align}

\paragraph{Momentum in $z$ (Plan (3), $z$-component).}
\begin{align}
M_{21}&=-ik\eta\,\mathcal G\,(c_{13}+c_{44}), &
M_{22}&=-\mathcal N\rho\omega^2+\mathcal G\,(c_{44}k^2-c_{33}\eta^2),\\
M_{23}&=e_{15}k^2-e_{33}\eta^2, &
M_{25}&=\eta\,\beta_3 .
\end{align}

\paragraph{Gauss law (Plan (4)).}
\begin{align}
M_{31}&=ik\eta\,(e_{15}+e_{31}), &
M_{32}&=e_{33}\eta^2-e_{15}k^2, &
M_{33}&=\epsilon_{11}k^2-\epsilon_{33}\eta^2,\\
M_{34}&=q, &
M_{35}&=\eta\,p_3 .
\end{align}

\paragraph{Carrier continuity (Plan (5)).}
\begin{align}
M_{43}&=b_nn_0\,L, &
M_{44}&=-i\omega-D_nL-ik\,v_d+\frac1{\tau_R}, &
M_{45}&=-\beta_{\rm th}L .
\end{align}

\paragraph{Heat (Plan (6)).}
\begin{align}
M_{51}&=f\,T_0\,\omega k\,\beta_1, &
M_{52}&=-i\omega f\,T_0\,\beta_3\,\eta,\\
M_{53}&=i\omega f\,T_0\,p_3\,\eta, &
M_{54}&=f\!\left(-i\omega\,q\,\Pi_{\rm th}+\frac{E_g}{\tau_R}\right),\\
M_{55}&=-i\omega f\rho c_p-k_{\rm th}L .
\end{align}

\section{The characteristic equation and what it says about the roots}

Non-trivial amplitudes exist only if
\begin{equation}
\det\mathbf M(\eta)=0 .
\label{eq:char}
\end{equation}
This is a polynomial in $\eta$. With $l_g\neq0$ its degree is $14$. With $l_g=0$ (nonlocal operator alone)
it drops to $10$; the gradient operator is what raises the order, the nonlocal one acting on the inertia does not.
Without the size effect the degree is $10$ as well. Please count the degree yourself from the entries above.

Two statements about the roots are given here without proof, and you should test both.
First, when $\beta_{\rm th}=\Pi_{\rm th}=0$, the polynomial is even in $\eta$: the roots come in pairs $\pm\eta$.
Second, when the thermodiffusion and Peltier coefficients are non-zero together with the piezoelectric or
pyroelectric constants, the polynomial has odd powers of $\eta$ as well, the roots no longer come in pairs, and
the plate modes cannot be split into symmetric and antisymmetric families. The reason is that a polar crystal
has no mirror symmetry about $z=0$ once the sign of the carrier charge matters. Because of this, the
determinant of the boundary matrix below is taken over all $14$ (or $10$) roots at once, without splitting it.
Switching $\beta_{\rm th},\Pi_{\rm th}$ off restores the pairs and the split.

\section{The solution through the thickness}

Let $\eta_1,\dots,\eta_{14}$ be the roots of Eq.~\eqref{eq:char}. For each root, the amplitude vector
$\mathbf a_m=(X_m,Z_m,\Phi_m,\nu_m,\Theta_m)^T$ is any non-zero column of the adjugate of
$\mathbf M(\eta_m)$, that is, $a_{m,j}=C_{\alpha j}(\eta_m)$ for one fixed row index $\alpha$ of the cofactor matrix
$C$. The general solution is
\begin{equation}
\begin{pmatrix}u_x\\u_z\\\varphi\\\Delta n\\\theta\end{pmatrix}
=\sum_{m=1}^{14}A_m\begin{pmatrix}X_m\\Z_m\\\Phi_m\\\nu_m\\\Theta_m\end{pmatrix}
e^{\eta_mz}\,e^{i(kx-\omega t)} ,
\label{eq:solution}
\end{equation}
with fourteen unknown constants $A_m$. When $l_g=0$ the sum has ten terms. Source: this is the standard
partial-wave construction for a layer; the general form is in every text on guided waves in plates. The
particular reduction to $\mathbf M(\eta)$ is ours.

\section{Boundary conditions and the final matrix}

\subsection{The seven conditions on each face}

Seven conditions are applied at each of $z=+h$ and $z=-h$, fourteen in all, one for each unknown $A_m$.
All are homogeneous. For a root $\eta_m$ write $\mathcal N_m=\mathcal N(\eta_m)$, $\mathcal G_m=\mathcal G(\eta_m)$.
The seven row functions are
\begin{align}
b_1&=\frac1{\mathcal N_m}\Bigl[\hat\sigma_{zz}-l_g^2\,ik\eta\,c_{44}(\eta X_m+ikZ_m)\Bigr],\label{eq:b1}\\
b_2&=\frac1{\mathcal N_m}\Bigl[\hat\sigma_{xz}-l_g^2\,ik\eta\,(ik\,c_{11}X_m+\eta\,c_{13}Z_m)\Bigr],\label{eq:b2}\\
b_3&=\frac{l_g^2}{\mathcal N_m}\,\eta\,c_{44}(\eta X_m+ikZ_m),\label{eq:b3}\\
b_4&=\frac{l_g^2}{\mathcal N_m}\,\eta\,(ik\,c_{13}X_m+\eta\,c_{33}Z_m),\label{eq:b4}\\
b_5&=ik\,e_{31}X_m+\eta\,e_{33}Z_m-\eta\,\epsilon_{33}\Phi_m+p_3\Theta_m,\label{eq:b5}\\
b_6&=-D_n\eta\nu_m+b_nn_0\,\eta\,\Phi_m-\beta_{\rm th}\,\eta\,\Theta_m,\label{eq:b6}\\
b_7&=\eta\,\Theta_m ,\label{eq:b7}
\end{align}
where in $b_1$ and $b_2$ the stresses $\hat\sigma_{zz},\hat\sigma_{xz}$ are those of Section~3 evaluated at
$(X_m,Z_m,\Phi_m,\Theta_m)$ and $\eta_m$, and $\eta=\eta_m$ throughout. Their physical meaning and sources are the
following.

Rows $b_1$ and $b_2$ are the normal and shear tractions on a free face, including the surface-gradient
terms of a strain-gradient continuum, and with the nonlocal operator inverted root by root, since
$\mathcal N^{-1}$ acts on $e^{\eta z}$ as $1/\mathcal N(\eta)$. For the traction of a strain-gradient continuum,
including the term with the tangential derivative, see \cite[Sec.~2, Eq.~(14)]{Charal2019}; the original is
\cite{Mindlin1964} (equation not read by us). For the nonlocal differential law and the constitutive boundary
conditions that go with it see \cite[Eqs.\ (18)--(19)]{Barretta2018} (read by us) and \cite{Lim2015}
(not seen). This pairing of the two papers for a free plate is a modelling choice made by us. It is the
one point where a different, equally defensible choice exists, and you will have to defend it in the paper.

Rows $b_3$ and $b_4$ are the double tractions, the quantities the strain-gradient term does work on at the
face: $R=n\cdot\mu\cdot n$ \cite[Sec.~2, condition b(i)]{Charal2019}. They exist only when $l_g\ne0$.

Row $b_5$: the normal electric displacement vanishes, an electrically open face (we neglect the field in the
vacuum outside; the other common choice, a shorted face, replaces $b_5$ by $\hat\varphi=0$). Ours; it is the
usual open-circuit condition for a piezoelectric plate.

Row $b_6$: no carrier flux through the face. The particle flux in the $z$-direction is
$\Gamma_z=-D_n\partial_z\Delta n-b_nn_0E_z-\beta_{\rm th}\partial_z\theta$, which is the drift--diffusion current of
\cite[Eq.\ (3), Eq.\ (7)]{Zhu2023} divided by the charge, with the thermodiffusion term of Plan (5) added.
Setting $\Gamma_z=0$ is ours.

Row $b_7$: adiabatic face, $\partial_z\theta=0$. Because $\Gamma_z=0$ as well, the Peltier part of the heat
current vanishes at the face and no cross term remains. For an isothermal face replace $b_7$ by $\Theta_m$.

\subsection{The final matrix}

Assemble the $14\times14$ matrix $\mathbf B$ with one row per condition and face and one column per root,
\begin{equation}
\mathbf B(k,\omega)=
\begin{pmatrix}
b_r(\eta_m)\,e^{+\eta_mh}\\[2pt]
b_r(\eta_m)\,e^{-\eta_mh}
\end{pmatrix}_{\;r=1..7,\ m=1..14},
\qquad \mathbf B\,\mathbf A=\mathbf 0 ,\quad \mathbf A=(A_1,\dots,A_{14})^T .
\label{eq:Bmatrix}
\end{equation}
The upper seven rows belong to $z=+h$ and the lower seven rows to $z=-h$. The dispersion relation is
\begin{equation}
\det\mathbf B(k,\omega)=0 ,
\label{eq:disp}
\end{equation}
to be solved for complex $k$ at real $\omega$. Attenuation per unit length is $k_i$. When $l_g=0$ use rows
$b_1,b_2,b_5,b_6,b_7$ and ten roots, giving a $10\times10$ matrix. The columns belonging to roots with
very large $|\mathrm{Re}\,\eta_m|$ (the thermal and carrier boundary layers are of this kind) overflow in
floating point; scaling each column by $e^{-|\mathrm{Re}\,\eta_m|h}$ leaves the zeros of the determinant
unchanged and avoids it.

\section{What to do with this file}

Go through Sections 3 to 7 in order and rederive each formula from the plan. Write a short list of every difference
you find, with the line of the source that supports your version. If you find none, say so, and say which of the
tests below you ran. Suggested tests, none of which needs a number from us:
(i) with all couplings off and $l_g=l_{\rm nl}=0$, $\det\mathbf B=0$ must reproduce the Rayleigh--Lamb relation
of an isotropic plate; (ii) with only the piezoelectric terms on, the thin-plate limit must give the
open-circuit stiffened bulk speed; (iii) with the carrier terms off and the heat equation in Fourier form
($\tau_q=0$), a plane longitudinal wave in the bulk must show thermoelastic damping with the sign and
magnitude you get from the standard formula; (iv) with only the carrier terms and the piezoelectric shear
coupling on, the bulk attenuation must peak at a carrier density that you can predict from the relaxation of the
free charge, $\sigma/\epsilon\approx\omega$; (v) letting $l_g,l_{\rm nl}\to0$ must reproduce the
classical result continuously.

\begin{thebibliography}{99}
\small
\bibitem{Zhu2023} J.~Zhu, M.~Negahban, J.~Xu, R.~Xia, Z.~Li, \emph{Micromachines} \textbf{14} (2023) 2174, \doi{10.3390/mi14122174}. Section 2.1, Eqs.\ (1)--(7).
\bibitem{Quintanilla2019} R.~Quintanilla, \emph{Math.\ Mech.\ Solids} \textbf{24} (2019) 4020, \doi{10.1177/1081286519862007}. Eqs.\ (2.5), (3.1)--(3.5).
\bibitem{Charal2019} A.~Charalambopoulos, E.~Douka, M.~Mavratzas, Second order formulation of boundary value problems in gradient elasticity, arXiv:1909.10776 (2019), Sec.~2, Eq.~(14).
\bibitem{Barretta2018} R.~Barretta, F.~Marotti de Sciarra, \emph{Int.\ J.\ Eng.\ Sci.} \textbf{130} (2018) 187, \doi{10.1016/j.ijengsci.2018.05.009}.
\bibitem{Lim2015} C.~W.~Lim, G.~Zhang, J.~N.~Reddy, \emph{J.\ Mech.\ Phys.\ Solids} \textbf{78} (2015) 298, \doi{10.1016/j.jmps.2015.02.001}.
\bibitem{Mindlin1964} R.~D.~Mindlin, \emph{Arch.\ Rational Mech.\ Anal.} \textbf{16} (1964) 51.
\end{thebibliography}
\end{document}
